\chapter{Verwandte Arbeiten}
\label{ch:relatedwork}

Dieses Kapitel ordnet \texttt{claude-tunnel} in bestehende Systeme ein und
schärft die zentrale Unterscheidung der Arbeit heraus: die
\emph{Providerblindheit}. Betrachtet werden vier Familien verwandter Systeme ---
klassische VPNs (Abschnitt~\ref{sec:rw-vpn}), Reverse-Tunnel-Dienste
(Abschnitt~\ref{sec:rw-reverse}), Anonymitäts- und Zensurumgehungssysteme
(Abschnitt~\ref{sec:rw-anon}) sowie QUIC-basierte Relays
(Abschnitt~\ref{sec:rw-quic}). Der abschließende Abschnitt~\ref{sec:rw-abgrenzung}
verdichtet die Abgrenzung in einer Gegenüberstellung. Leitfrage ist durchgehend:
\emph{Wer terminiert die Verschlüsselung, und wer kann folglich Klartext lesen
oder die Verbindung kappen?}

\section{Klassische VPNs und WireGuard}
\label{sec:rw-vpn}

Klassische VPNs stellen einen verschlüsselten Tunnel zwischen zwei Endpunkten
her. Bei WireGuard~\cite{donenfeld2017wireguard} etwa terminiert die Verschlüsselung an den beiden Peers
selbst: Der Tunnel verbindet einen Client mit einem Server, der zugleich
Gegenstelle und Vertrauensanker ist. Für den hier verfolgten Anwendungsfall ---
ein kundeneigenes Origin hinter einer restriktiven Netzgrenze ohne öffentliche
Adresse erreichbar zu machen --- ist dieses Modell nur bedingt geeignet: Es
existiert kein providerblindes \emph{Relay} zwischen den Parteien, das Verkehr
vermittelt, ohne selbst Endpunkt zu sein. Wer einen VPN-Server als
Vermittlungsknoten betreibt, terminiert dessen Verschlüsselung und sieht den
weitergeleiteten Verkehr im Klartext. WireGuard löst die Transport- und
Schlüsselfrage elegant, adressiert aber gerade nicht die Trennung von blinder
Vermittlung und kundenseitig verankerter Ende-zu-Ende-Sicherheit, die den Kern
dieser Arbeit bildet. \texttt{claude-tunnel} übernimmt die Idee moderner,
schlanker Handshakes --- die Datenebene nutzt Noise~\cite{noiseprotocol} ---,
verschiebt die Terminierung jedoch vom Vermittler an das kundeneigene Ziel.

Näher am hier verfolgten Ansatz als das klassische VPN liegen jüngere Overlay-
und Secure-Channel-Systeme, deren Vermittler den Nutzlast-Schlüssel bewusst nicht
halten. Nebula~\cite{nebula} spannt ein Noise-basiertes Mesh zwischen
authentifizierten Hosts auf, in dem koordinierende \emph{Lighthouses} die Peers
zusammenführen, ohne die Ende-zu-Ende-Sitzung zu terminieren. Tailscale reicht
Verkehr, wenn kein direkter Pfad zustande kommt, über verschlüsselte
\emph{DERP}-Relays weiter, die die WireGuard-Pakete lediglich
durchreichen~\cite{tailscalederp}; die zugehörige Steuerebene lässt sich mit
Headscale~\cite{headscale} selbst betreiben. Ockam~\cite{ockam} wiederum baut
mehrsprüngige, gegenseitig authentifizierte Ende-zu-Ende-Kanäle, deren Relays
als Architektur-Eigenschaft ausschließlich Chiffretext weiterleiten. Diese
Systeme zeigen, dass \emph{payload-blindes Relaying} für sich genommen kein
Alleinstellungsmerkmal ist, sondern verbreitete Praxis. Der Beitrag dieser Arbeit
liegt daher nicht in der Nutzlast-Blindheit allein, sondern in ihrer Kombination
mit einem opaken Routing-Token (das zusätzlich das Vermittlungsziel verbirgt),
einem KYC-freien Proof-of-Work-Rendezvous und kundenverankerten Noise-Schlüsseln
ohne zentrale PKI (Abschnitt~\ref{sec:rw-abgrenzung}).

\section{Reverse-Tunnel-Dienste}
\label{sec:rw-reverse}

Reverse-Tunnel-Dienste wie Cloudflare Tunnel~\cite{cloudflaretunnel},
ngrok~\cite{ngrok} oder Tailscale Funnel~\cite{tailscalefunnel} lösen
dieselbe Erreichbarkeitsaufgabe: Ein Agent im privaten Netz baut eine
ausgehende Verbindung zu einer Anbieter-Infrastruktur auf, über die eingehender
Verkehr an das interne Ziel zurückgereicht wird. Damit entfällt die
Notwendigkeit eingehender Firewall-Freigaben. Entscheidend für die Abgrenzung ist
jedoch, \emph{wo die Transport\-verschlüsselung terminiert}: Bei diesen Diensten
endet die TLS-Sitzung beim Anbieter. Der Betreiber steht als vollwertiger
Man-in-the-Middle im Pfad, sieht den Klartext des vermittelten Verkehrs und kann
einzelne Tunnel jederzeit inhaltlich inspizieren oder kappen. Adressiert werden
die Tunnel überdies typischerweise über öffentliche Hostnamen, deren Zuordnung
der Anbieter kennt und kontrolliert.

Genau dieses Vertrauensmodell ist das \emph{Gegenmodell} zu \texttt{claude-tunnel}.
Der providerblinde Ansatz kehrt die Eigenschaft um: Das Relay leitet
ausschließlich Chiffretext weiter, hält keinen Schlüssel zur Nutzlast und
adressiert Tunnel nicht über einen Hostnamen, sondern über ein opakes
\emph{Routing-Token} ohne SNI-Information. Der Betreiber bleibt damit von Inhalt
und Vermittlungsziel abgeschnitten und ist als Druck- oder Zensurpunkt entwertet
--- die funktionale Bequemlichkeit des Reverse-Tunnels wird ohne die
Vertrauensbürde des klartextsichtigen Anbieters erreicht.

\section{Anonymität und Zensurumgehung}
\label{sec:rw-anon}

Tor~\cite{dingledine2004tor} und allgemeiner das Onion-Routing verfolgen ein anderes Bedrohungsmodell:
Ziel ist die \emph{Anonymität} der kommunizierenden Parteien, erreicht durch
mehrstufige, geschichtete Verschlüsselung über mehrere unabhängige Relays, sodass
kein einzelner Knoten Sender und Empfänger zugleich kennt. \emph{Pluggable
transports} wie obfs4~\cite{obfs4} oder Shadowsocks~\cite{shadowsocks} ergänzen dies um Verschleierung des
Verkehrsmusters, um Deep-Packet-Inspection und Protokollfilter zu unterlaufen.
Weitere Umgehungsansätze verlagern die Resistenz in die Netzinfrastruktur ---
etwa Domain Fronting~\cite{fifield2015domainfronting}, das Sperren durch
Mitnutzung großer, nicht blockierbarer Fronting-Domains unterläuft, oder
Refraction-/Decoy-Routing wie Telex~\cite{wustrow2011telex}; Messplattformen wie
OONI~\cite{filasto2012ooni} dokumentieren Verbreitung und Methodik solcher
Blockaden empirisch. Diese Systeme schützen also primär die \emph{Identität} und
die \emph{Beobachtbarkeit} der Kommunikation gegenüber Netzbeobachtern.

\texttt{claude-tunnel} zielt demgegenüber nicht auf Anonymität, sondern auf
\emph{Providerblindheit}: Der Betreiber darf durchaus wissen, \emph{dass} ein
Tunnel existiert und grobe Zeit- und Volumenzähler sehen --- er darf nur die
Nutzlast weder lesen noch das Origin-Schlüsselmaterial erlangen. Die
Bedrohungsmodelle sind damit komplementär, nicht deckungsgleich. Verwandt bleiben
gleichwohl die Fragen der \emph{Metadaten}-Minimierung (welche strukturellen
Spuren ein Vermittler zwangsläufig sieht) und der \emph{Denial-of-Service}-
Resistenz eines offenen, KYC-freien Zugangs. Wo Tor auf Netzwerk- und
Reputationsmechanismen setzt, verwendet \texttt{claude-tunnel} ein
Proof-of-Work-Rendezvous~\cite{back2002hashcash} mit tokenweiser Ratenbegrenzung
als primären Sybil-~\cite{douceur2002sybil} und Flutungs-Schutz.

\section{QUIC-Relays und MASQUE}
\label{sec:rw-quic}

Auf der Transportebene ist \texttt{claude-tunnel} den QUIC-basierten
Relay-Ansätzen verwandt. QUIC~\cite{rfc9000} bündelt Verbindungsaufbau,
Verschlüsselung und Multiplexing über UDP und eignet sich damit für den
schnellen Aufbau vieler kurzlebiger Tunnel; die Vertraulichkeit ruht dabei auf
TLS~1.3~\cite{rfc8446}, das QUIC in seinen Handshake integriert. Das
MASQUE-Framework~\cite{rfc9298,rfc9484} baut auf QUIC
und HTTP/3 auf, um Proxy- und Tunneldienste über eine QUIC-Verbindung zu
transportieren --- konzeptionell ein Relay, das UDP- oder IP-Verkehr für einen
Client weiterreicht. Die Datenebene dieser Arbeit teilt diese Transportbasis:
QUIC (via \texttt{quinn}) ist der primäre Transport, ergänzt um einen
TLS-über-TCP-Rückfall für Umgebungen, in denen UDP blockiert ist.

Der Unterschied liegt erneut in der Terminierung und der Schlüsselverankerung:
Ein MASQUE-Proxy terminiert die QUIC-Verbindung des Clients bei sich und ist
somit prinzipiell klartextsichtig auf der Proxy-Sitzung. \texttt{claude-tunnel}
legt über den QUIC-Transport eine zweite, providerblinde Schicht: eine
Ende-zu-Ende-Noise-Sitzung zwischen Client und Origin, deren Chiffretext das
Relay lediglich durchreicht. Der Transport ist verwandt; das Vertrauensmodell ist
es nicht.

\section{Abgrenzung}
\label{sec:rw-abgrenzung}

Tabelle~\ref{tab:rw-abgrenzung} stellt die vier Familien anhand der
entscheidenden Merkmale gegenüber. Der Unterschied bündelt sich in einer Zeile:
Bei WireGuard, Cloudflare Tunnel und Tor liegt die für den Nutzlastschutz
relevante Terminierung entweder beim Vermittler oder ist gar nicht als blindes
Relay konzipiert; nur \texttt{claude-tunnel} verbindet ein nutzlast-blindes Relay
mit einer kundenseitig am Origin verankerten Terminierung.

\begin{table}[t]
\centering
\caption{Abgrenzung entlang von Terminierung, Klartextsicht, Kappbarkeit und
Adressierung.}
\label{tab:rw-abgrenzung}
\begin{tabular}{llll}
\toprule
Eigenschaft & WireGuard & Cloudflare Tunnel & Tor \\
\midrule
Wer terminiert TLS/Krypto & Peer/Server & Anbieter & Origin (mehrfach) \\
Wer sieht Klartext & Server-Peer & Anbieter & kein Einzelknoten \\
Wer kann kappen & Server-Peer & Anbieter & jedes Relay (Pfad) \\
Öffentl. Hostname nötig & Server-Endpunkt & ja & nein \\
\bottomrule
\end{tabular}
\end{table}

Die konzeptionell nächsten Vorläufer sind nicht die klassischen Tunnel, sondern
jüngere Ansätze, die das \emph{Wissen} über eine Verbindung bewusst auf mehrere
Parteien aufteilen. Oblivious HTTP~\cite{rfc9458} trennt den Anfrageinhalt von
der Absenderidentität: Ein Relay verbirgt die Client-Adresse, während ein Gateway
die Anfrage entschlüsselt und an das Ziel weiterreicht. Damit sieht der Gateway
sehr wohl Ziel und Klartext, und das Modell ist auf einzelne, zustandslose
HTTP-Anfragen zugeschnitten statt auf einen langlebigen, protokollagnostischen
Tunnel. Apples iCloud Private Relay~\cite{appleprivaterelay} nutzt zwei Hops mit
Ingress-/Egress-Trennung, sodass kein einzelner Betreiber zugleich Client-IP und
Zieladresse kennt; die Trennung ruht jedoch auf \emph{zwei getrennten
Betreiberrollen}, und die Sitzung zum Ziel terminiert regulär --- eine
kundenseitig verankerte Nutzlast-Blindheit gegenüber einem \emph{einzelnen}
Vermittler ist damit nicht gegeben. \texttt{claude-tunnel} grenzt sich gegen
beide dadurch ab, dass ein \emph{einzelner} Betreiber weder Nutzlast noch
Vermittlungsziel erfährt --- Chiffretext statt Klartext, opakes Routing-Token
statt Hostname --- und die Terminierung unabhängig von einer
Zwei-Parteien-Aufteilung am kundeneigenen Origin verankert bleibt.

Die Alleinstellung von \texttt{claude-tunnel} ergibt sich aus der Kombination
vierer Bausteine, die in dieser Zusammenstellung bei keinem der verglichenen
Systeme --- auch nicht bei diesen nächsten Vorläufern --- auftritt:

\begin{itemize}
  \item ein \emph{nutzlast-blindes Relay}, das ausschließlich Chiffretext
    weiterleitet und keinen Schlüssel zur Nutzlast besitzt;
  \item ein \emph{opakes Routing-Token} statt eines Hostnamens, das dem
    Betreiber auch das Vermittlungsziel verbirgt (keine SNI- oder
    Hostnamen-Information);
  \item ein \emph{Proof-of-Work-Rendezvous} mit tokenweiser Ratenbegrenzung als
    KYC-freier Zugangs- und DoS-Schutz~\cite{back2002hashcash};
  \item \emph{kundenseitig verankerte Noise-Schlüssel ohne
    PKI}~\cite{noiseprotocol} --- die Ende-zu-Ende-Sicherheit terminiert am
    kundeneigenen Origin und benötigt weder eine Zertifikatshierarchie noch einen
    anbieterseitigen Vertrauensanker.
\end{itemize}

Für \texttt{claude-tunnel} bedeutet dies: Die funktionale Erreichbarkeit eines
Reverse-Tunnels wird mit dem Vertrauensmodell eines providerblinden Relays
zusammengeführt --- der Betreiber vermittelt, ohne Endpunkt oder Vertrauensanker
zu sein.
